% संपूर्ण मराठी ग्रंथातील प्रकरण 37: चर्च–रॉसर गुणधर्म
% हा संपादनयोग्य स्रोतखंड आहे. संकलनासाठी संपूर्ण स्रोतसंग्रहातील
% अक्षररूपे, पूर्वभाग, चिन्हव्याख्या आणि आकृती-स्रोत आवश्यक आहेत.
% मूळ: Open Logic Project; मजकूर आणि रूपांतर CC BY 4.0.
% यंत्रानुवाद, दुरुस्ती आणि तपासणी: OpenAI Codex —
% GPT-5.6 Sol आणि GPT-6 Sol, दोन्ही Ultra effort.
% दुरुस्ती-पुनरावलोकन: GPT-6.1 Sol, Ultra effort.
\chapter{चर्च–रॉसर गुणधर्म}\label{lam:cr:chap}









\section{व्याख्या आणि गुणधर्म}\label{lam:cr:dap:sec}

या प्रकरणात चर्च–रॉसर गुणधर्माची कल्पना आणि त्या
गुणधर्माचे काही सामान्य पैलू मांडू.

\begin{defn}[चर्च–रॉसर गुणधर्म, CR] पदांवरील संबंध $\xredone$
  पुढील अट पूर्ण करत असेल, तर त्याला \emph{चर्च–रॉसर गुणधर्म}
  आहे असे म्हणतात: $M \xredone P$ आणि $M \xredone Q$
  असतील, तर असे~$N$ अस्तित्वात असते की
  $P \xredone N$ आणि $Q \xredone N$.
\end{defn}

लॅम्डा कलनाकडे संगणनाचे एक प्रतिमान म्हणून पाहता येते.
त्यात सामान्य रूपातील पदे ``मूल्ये'' आहेत आणि पदांवरील
न्यूनीकरणसंबंध हे ``संगणननियम'' आहेत. चर्च–रॉसर गुणधर्म
असे सांगतो की संगणन पुढे नेण्याचे एकाहून अधिक मार्ग
असले, तरी निकाल मिळत असल्यास त्याचे मूल्य एकमेव असते.

प्राथमिक बीजगणितातील उदाहरण घ्या. $4 \times (1+2) + 3$
मोजण्याचे अनेक मार्ग आहेत. आधी $1+2$ चे~$3$ करून
$4
\times 3+3$ मिळवता येते. किंवा वितरणनियम वापरून
आधी $4 \times (1+2)$ चे $4 \times 1+4
\times 2+3$ असे रूप करता येते. या दोन्ही रूपांचे
पुढील न्यूनीकरण $12+3$ मध्ये होते.

$\xredone$ म्हणजे $\beta$-न्यूनीकरण घेतले, तर
चर्च–रॉसर गुणधर्माचा एक परिणाम लगेच दिसतो:
पदाचे सामान्य रूप असेल, तर ते एकमेव असते.
समजा $M$ चे $P$ आणि $Q$ मध्ये न्यूनीकरण होते,
आणि ही दोन्ही सामान्य रूपे आहेत. चर्च–रॉसर गुणधर्मामुळे
असे $N$ असते की $P$ आणि $Q$ या दोन्हींचे त्यात
न्यूनीकरण होते. पण गृहीतकानुसार $P$ आणि $Q$ ही
सामान्य रूपे असल्याने $P$ आणि $Q$ यांचे $N$ मध्ये
न्यूनीकरण फक्त तुच्छ न्यूनीकरणच असू शकते.
म्हणजे $P$, $Q$ आणि $N$ ही समान आहेत.
त्यामुळे पदाच्या सामान्य रूपाचा उल्लेख आपण
\emph{एकमेव} म्हणून करू शकतो.

लॅम्डा कलनाकडे संगणनाचे प्रतिमान म्हणून पाहिल्यास,
एखाद्या पदाचे सामान्य रूप हा त्या पदापासून सुरू झालेल्या
संगणनाचा ``अंतिम निकाल'' मानता येतो. वरील निष्कर्षानुसार
असा अंतिम निकाल अस्तित्वात असेल, तर तो एकच असतो.
उदाहरणार्थ, $4 \times (1+2)+3$ मोजल्याचा निकाल
एकच, म्हणजे~$15$, असतो.

\begin{thm} \label{lam:cr:dap:thm:str}
  संबंध $\xredone$ ला चर्च–रॉसर गुणधर्म असेल आणि
  $\xred$ हा $\xredone$ समाविष्ट करणारा सर्वांत लहान
  संक्रमक संबंध असेल, तर $\xred$ लाही चर्च–रॉसर
  गुणधर्म असतो.
\end{thm}

\begin{proof}
  समजा
  \begin{align*}
    M & \xredone P_1 \xredone \dots \xredone P_m \text{ आणि}\\
    M & \xredone Q_1 \xredone \dots \xredone Q_n.
  \end{align*}
  उंची $m + 1$ आणि रुंदी $n + 1$ असलेली पदांची
  $N$ जाळी रचून प्रमेय सिद्ध करू. $i$ व्या ओळीतील
  आणि $j$ व्या स्तंभातील पद $N_{i,j}$ असे दर्शवू.

  $N_{i,j} \xredone N_{i+1,j}$ आणि
  $N_{i,j} \xredone N_{i,j+1}$ होईल अशी $N$
  रचू. तिची व्याख्या पुढीलप्रमाणे:
  \begin{align*}
    N_{0,0} &= M \\
    N_{i,0} &= P_i && \text{जर } 1 \le i \le m \\
    N_{0,j} &= Q_j && \text{जर } 1 \le j \le n \\
  \intertext{आणि इतर प्रसंगी:}
    N_{i,j} &= R
  \end{align*}
  येथे $R$ हे असे पद आहे की $N_{i-1,j} \xredone R$
  आणि $N_{i,j-1}
  \xredone R$. $\xredone$ च्या चर्च–रॉसर गुणधर्मामुळे
  असे पद नेहमी अस्तित्वात असते.

  आता $N_{m,0} \xredone \dots \xredone N_{m,n}$
  आणि $N_{0,n}
  \xredone \dots \xredone N_{m,n}$ मिळते.
  येथे $N_{m,0}$ हे $P_m$ आणि $N_{0,n}$ हे $Q_n$ आहे.
  $\xred$ च्या व्याख्येनुसार प्रमेय सिद्ध होते.
\end{proof}












\section{समांतर $\beta$-न्यूनीकरण}\label{lam:cr:pb:sec}

या विभागात \emph{समांतर $\beta$-न्यूनीकरण} ही संकल्पना
मांडू आणि तिला चर्च–रॉसर गुणधर्म आहे हे सिद्ध करू.

\begin{defn}[समांतर $\beta$-न्यूनीकरण, $\bredpar$] \label{lam:cr:pb:defn:bredpar}
  पदांचे समांतर न्यूनीकरण ($\bredpar$) पुढीलप्रमाणे
  विगमनाने परिभाषित केले जाते:
  \begin{enumerate}
    \item \label{lam:cr:pb:defn:bredpar1} $x \bredpar x$.
    \item \label{lam:cr:pb:defn:bredpar2} $N \bredpar N'$ असेल, तर $\lambd[x][N] \bredpar
      \lambd[x][N']$.
    \item \label{lam:cr:pb:defn:bredpar3} $P \bredpar P'$ आणि $Q \bredpar Q'$ असतील, तर $PQ \bredpar
      P'Q'$.
    \item \label{lam:cr:pb:defn:bredpar4} $N \bredpar N'$ आणि $Q \bredpar Q'$ असतील, तर
      $(\lambd[x][N])Q \bredpar \Subst{N'}{Q'}{x}$.
  \end{enumerate}
\end{defn}

समांतर $\beta$-न्यूनीकरणात एका पायरीत पदातील कितीही
रेडेक्सचे संकुचन करता येते. ते $\beta$-न्यूनीकरणापेक्षा
वेगळे आहे: मूळ पदात असलेल्या रेडेक्सचेच संकुचन त्या
पायरीत करता येते; समांतर $\beta$-न्यूनीकरणामुळे नव्याने निर्माण
झालेल्या रेडेक्सचे नाही. उदाहरणार्थ,
$(\lambd[f][fx])(\lambd[y][y])$ या पदाचे समांतर
$\beta$-न्यूनीकरण केल्यास तेच पद किंवा
$(\lambd[y][y])x$ मिळू शकते; एका पायरीत पुढे~$x$
मिळू शकत नाही. मात्र नेहमीच्या $\beta$-न्यूनीकरणाने
ते~$x$ मध्ये न्यूनीत होते. कारण संबंधित रेडेक्स समांतर
$\beta$-न्यूनीकरणाची पहिली पायरी झाल्यावरच निर्माण
होतो. दुसरी समांतर $\beta$-न्यूनीकरणाची पायरी
घेतल्यास~$x$ मिळते.

\begin{thm}\label{lam:cr:pb:thm:refl}
  $M \bredpar M$.
\end{thm}

\begin{proof}
  सरावासाठी.
\end{proof}

\begin{prob}
  \ref{lam:cr:pb:thm:refl} सिद्ध करा.
\end{prob}

\begin{defn}[$\beta$-पूर्ण विकास]\label{lam:cr:pb:defn:bcd}
  $M$ चा \emph{$\beta$-पूर्ण विकास}~$\bcd{M}$ पुढीलप्रमाणे
  विगमनाने परिभाषित केला जातो:
  \begin{align}
    \bcd{x} &= x \label{lam:cr:pb:defn:bcd1} \\
    \bcd{(\lambd[x][N])} &= \lambd[x][\bcd{N}] \label{lam:cr:pb:defn:bcd2}\\
    \bcd{(PQ)} &= \bcd{P}\bcd{Q} && \text{जर $P$ हे $\lambd$-अमूर्तीकरण नसेल तर}
    \label{lam:cr:pb:defn:bcd3} \\
    \bcd{((\lambd[x][N])Q)} &= \Subst{\bcd{N}}{\bcd{Q}}{x} \label{lam:cr:pb:defn:bcd4}
  \end{align}
\end{defn}

पदाचा $\beta$-पूर्ण विकास म्हणजे नावाप्रमाणे
``संपूर्ण समांतर न्यूनीकरण''. समांतर $\beta$-न्यूनीकरणात
एखाद्या रेडेक्सचे संकुचन टाळता येते; पण पूर्ण विकासात
सर्व मूळ रेडेक्सचे संकुचन करावेच लागते. म्हणून
$(\lambd[f][fx])(\lambd[y][y])$ चा पूर्ण विकास तेच
पद नसून $(\lambd[y][y])x$ आहे.

\begin{editorial}
  या व्याख्येत अडचण अशी की ($\lambd$-)पदांवरील फलने
  पुनरावर्ती रीतीने कशी परिभाषित करायची हे अजून
  मांडलेले नाही. पुढे दुरुस्त करू.
\end{editorial}

\begin{lem}\label{lam:cr:pb:lem:comp}
  $M \bredpar M'$ आणि $R \bredpar R'$ असतील, तर $\Subst{M}{R}{y}
  \bredpar \Subst{M'}{R'}{y}$.
\end{lem}

\begin{proof}

अमूर्तीकरणांची नावे आवश्यक त्या सर्व मुक्त चरांपासून
आणि आदेशित चरांपासून वेगळी निवडू. पुढील आदेशनसूत्र वापरू:
\[\Subst{\Subst{N}{Q}{x}}{R}{y}
=\Subst{\Subst{N}{R}{y}}{\Subst{Q}{R}{y}}{x},
\qquad x\ne y,\quad x\notin\FV{R}.
\]
हे सूत्र पदाच्या रचनेवरील विगमनाने सिद्ध होते. चराचे नाव
$x$ असल्यास दोन्ही बाजू $\Subst{Q}{R}{y}$ होतात;
नाव $y$ असल्यास दोन्ही बाजू $R$ होतात, कारण $R$ मध्ये
$x$ मुक्त नाही. इतर चर तसेच राहतात. उपयोजनात दोन्ही
उपपदांवर विगमन गृहीतक लागू होते. अमूर्तीकरणात त्याचे
बद्ध नाव दोन्ही आदेशित पदांच्या मुक्त चरांपासून आणि
$x,y$ पासून वेगळे निवडल्याने तेच गृहीतक देहाला लागू
करता येते. याने वर्गपदांवरील सूत्र सिद्ध होते.
या सिद्धतेतील बाह्य अमूर्तीकरणाचे नाव $x$ हे $y$ पासून
आणि $R,R'$ यांच्या सर्व मुक्त चरांपासून वेगळे निवडू.
  $M \bredpar M'$ च्या निष्पत्ती वर विगमन करू.
  \begin{enumerate}
    \item अखेरची पायरी \ref{lam:cr:pb:defn:bredpar1} आहे: सरावासाठी.
    \item अखेरची पायरी \ref{lam:cr:pb:defn:bredpar2} आहे: मग $M$ हे
      $\lambd[x][N]$ आणि $M'$ हे $\lambd[x][N']$ आहे,
      जेथे $N \bredpar N'$. आपल्याला
      $\Subst{(\lambd[x][N])}{R}{y} \bredpar
      \Subst{(\lambd[x][N'])}{R'}{y}$, म्हणजे
      $\lambd[x][\Subst{N}{R}{y}] \bredpar
      \lambd[x][\Subst{N'}{R'}{y}]$, हे सिद्ध करायचे आहे.
      ते \ref{lam:cr:pb:defn:bredpar2} आणि विगमन गृहीतकावरून
      लगेच मिळते.
    \item अखेरची पायरी \ref{lam:cr:pb:defn:bredpar3} आहे: सरावासाठी.
    \item अखेरची पायरी \ref{lam:cr:pb:defn:bredpar4} आहे: $M$ हे
      $(\lambd[x][N])Q$ आणि $M'$ हे $\Subst{N'}{Q'}{x}$ आहे.
      आपल्याला $\Subst{((\lambd[x][N])Q)}{R}{y} \bredpar
      \Subst{\Subst{N'}{Q'}{x}}{R'}{y}$, म्हणजे
      $(\lambd[x][\Subst{N}{R}{y}])\Subst{Q}{R}{y} \bredpar
      \Subst{\Subst{N'}{R'}{y}}{\Subst{Q'}{R'}{y}}{x}$,
      हे सिद्ध करायचे आहे. वरील आदेशनसूत्राने या दोन
      उजव्या बाजू समान आहेत. मग \ref{lam:cr:pb:defn:bredpar4}
      आणि दोन्ही पूर्वपक्षांना लागू केलेल्या विगमन गृहीतकावरून
      आवश्यक समांतर न्यूनीकरण मिळते.
  \end{enumerate}
\end{proof}

\begin{prob}
  \ref{lam:cr:pb:lem:comp} ची सिद्धता पूर्ण करा.
\end{prob}

\begin{lem}\label{lam:cr:pb:lem:cont}
  $M \bredpar M'$ असेल, तर $M' \bredpar \bcd{M}$.
\end{lem}
\begin{proof}
  $M \bredpar M'$ च्या निष्पत्ती वर विगमन करू.
  \begin{enumerate}
    \item अखेरचा नियम \ref{lam:cr:pb:defn:bredpar1} आहे: सरावासाठी.
    \item अखेरचा नियम \ref{lam:cr:pb:defn:bredpar2} आहे: $M$ हे
      $\lambd[x][N]$ आणि $M'$ हे $\lambd[x][N']$ आहे,
      जेथे $N \bredpar N'$. आपल्याला $\lambd[x][N'] \bredpar
      \bcd{(\lambd[x][N])}$, म्हणजे \ref{lam:cr:pb:defn:bcd2} नुसार
      $\lambd[x][N'] \bredpar \lambd[x][\bcd{N}]$,
      दाखवायचे आहे. ते \ref{lam:cr:pb:defn:bredpar2} आणि
      विगमन गृहीतकावरून मिळते.
    \item अखेरचा नियम \ref{lam:cr:pb:defn:bredpar3} आहे: $M$ हे $PQ$
      आणि $M'$ हे $P'Q'$ आहे, जेथे $P$, $Q$, $P'$, $Q'$
      अशी पदे आहेत की $P \bredpar P'$ आणि $Q \bredpar Q'$.
      विगमन गृहीतकानुसार $P' \bredpar \bcd{P}$ आणि
      $Q' \bredpar \bcd{Q}$.
      \begin{enumerate}
        \item काही $x$ आणि $N$ साठी $P$ हे $\lambd[x][N]$
          असेल, तर काही $N'$ साठी $P'$ हे $\lambd[x][N']$
          असावे लागते आणि $N \bredpar N'$. विगमन गृहीतकानुसार
          $N' \bredpar \bcd{N}$ आणि $Q' \bredpar \bcd{Q}$.
          मग \ref{lam:cr:pb:defn:bredpar4} नुसार
          $(\lambd[x][N'])Q' \bredpar
          \Subst{\bcd{N}}{\bcd{Q}}{x}$.
        \item $P$ हे $\lambd$-अमूर्तीकरण नसेल, तर
          \ref{lam:cr:pb:defn:bredpar3} नुसार $P'Q' \bredpar
          \bcd{P}\bcd{Q}$; आणि \ref{lam:cr:pb:defn:bcd3} नुसार
          उजवीकडील पद $\bcd{PQ}$ आहे.
      \end{enumerate}
    \item अखेरचा नियम \ref{lam:cr:pb:defn:bredpar4} आहे: काही
      $x$, $N$, $Q$, $N'$ आणि~$Q'$ साठी $M$ हे
      $(\lambd[x][N])Q$ आणि $M'$ हे $\Subst{N'}{Q'}{x}$,
      जेथे $N \bredpar N'$ आणि $Q \bredpar Q'$.
      विगमन गृहीतकानुसार $N' \bredpar \bcd{N}$ आणि
      $Q' \bredpar \bcd{Q}$. \ref{lam:cr:pb:lem:comp} नुसार
      $\Subst{N'}{Q'}{x} \bredpar \Subst{\bcd{N}}{\bcd{Q}}{x}$;
      उजवीकडील पद नेमके $\bcd{((\lambd[x][N])Q)}$ आहे.
  \end{enumerate}
\end{proof}

\begin{prob}
  \ref{lam:cr:pb:lem:cont} ची सिद्धता पूर्ण करा.
\end{prob}

\begin{thm}\label{lam:cr:pb:thm:cr}
  $\bredpar$ ला चर्च–रॉसर गुणधर्म आहे.
\end{thm}

\begin{proof}
  \ref{lam:cr:pb:lem:cont} वरून लगेच मिळते.
\end{proof}












\section{$\beta$-न्यूनीकरण}\label{lam:cr:b:sec}

\begin{lem}\label{lam:cr:b:lem:one-par}
  $M \bredone M'$ असेल, तर $M \bredpar M'$.
\end{lem}
\begin{proof} $M \bredone M'$ च्या निष्पत्तीवर विगमन करू.
  मूळस्थानी संकुचनाच्या प्रकरणात काही $x$, $N$ आणि~$Q$ साठी
  $M$ हे $(\lambd[x][N])Q$ आणि $M'$ हे
  $\Subst{N}{Q}{x}$ आहे. \ref{lam:cr:pb:thm:refl} नुसार
  $N \bredpar N$ आणि $Q \bredpar Q$ असल्याने
  \ref{lam:cr:pb:defn:bredpar}\ref{lam:cr:pb:defn:bredpar4} नुसार
  $(\lambd[x][N])Q \bredpar \Subst{N}{Q}{x}$ लगेच मिळते.
  अमूर्तीकरण किंवा उपयोजनाच्या आत संकुचन झाल्यास,
  विगमन गृहीतक, अपरिवर्तित उपपदाची समांतर परावर्तिता
  आणि संबंधित रचना-नियमांवरून निष्कर्ष मिळतो.
\end{proof}

\begin{lem}\label{lam:cr:b:lem:par-red}
  $M \bredpar M'$ असेल, तर $M \bred M'$.
\end{lem}

\begin{proof} $M \bredpar M'$ च्या निष्पत्ती वर विगमन करू.
  \begin{enumerate}
  \item अखेरचा नियम \ref{lam:cr:pb:defn:bredpar1} आहे: मग $M$
    आणि $M'$ दोन्ही $x$ आहेत आणि $x \bred x$.
  \item अखेरचा नियम \ref{lam:cr:pb:defn:bredpar2} आहे: काही
    $x$, $N$, $N'$ साठी $M$ हे $\lambd[x][N]$
    आणि $M'$ हे $\lambd[x][N']$ आहे, जेथे
    $N \bredpar N'$. विगमन गृहीतकानुसार $N \bred N'$.
    मग $N \bred N'$ साठी लागणारी तीच $\bredone$
    संकुचने अमूर्तीकरणाच्या आत करून
    $\lambd[x][N] \bred \lambd[x][N']$ मिळते.
  \item अखेरचा नियम \ref{lam:cr:pb:defn:bredpar3} आहे: काही
    $P$, $Q$, $P'$, $Q'$ साठी $M$ हे $PQ$ आणि
    $M'$ हे $P'Q'$ आहे, जेथे $P \bredpar P'$
    आणि $Q \bredpar Q'$. विगमन गृहीतकानुसार
    $P \bred P'$ आणि $Q \bred Q'$. म्हणून आधी
    $P \bred P'$ आणि नंतर $Q \bred Q'$ असे
    न्यूनीकरण केल्यावर $PQ \bred P'Q'$ मिळते.
  \item अखेरचा नियम \ref{lam:cr:pb:defn:bredpar4} आहे: काही
    $x$, $N$, $N'$, $Q$, $Q'$ साठी $M$ हे
    $(\lambd[x][N])Q$ आणि $M'$ हे
    $\Subst{N'}{Q'}{x}$ आहे, जेथे $N \bredpar N'$
    आणि $Q \bredpar Q'$. विगमन गृहीतकानुसार
    $Q \bred Q'$ आणि $N \bred N'$. म्हणून आधी
    $N \bred N'$, मग $Q \bred Q'$, आणि शेवटी
    $(\lambd[x][N'])Q'$ चे $\Subst{N'}{Q'}{x}$
    मध्ये संकुचन करून
    $(\lambd[x][N])Q \bred \Subst{N'}{Q'}{x}$ मिळते.
  \end{enumerate}
\end{proof}

\begin{lem}\label{lam:cr:b:lem:str}
  $\bredpar$ समाविष्ट करणारा सर्वांत लहान संक्रमक
  संबंध $\bred$ आहे.
\end{lem}

\begin{proof}
  $\bredpar$ समाविष्ट करणारा सर्वांत लहान संक्रमक
  संबंध $\xred$ मानू.

  $\bred \subseteq \xred$: समजा $M \bred M'$,
  म्हणजे $M \ident M_1 \bredone \dots \bredone M_k \ident M'$.
  \ref{lam:cr:b:lem:one-par} नुसार $M \ident M_1 \bredpar
  \dots \bredpar M_k \ident M'$. $\xred$ मध्ये
  $\bredpar$ समाविष्ट आहे आणि तो संक्रमक आहे;
  म्हणून $M \xred M'$. शून्य संकुचनांच्या प्रकरणातही
  समांतर न्यूनीकरणाच्या परावर्ती गुणधर्मामुळे हाच निष्कर्ष मिळतो.

  $\xred \subseteq \bred$: समजा $M \xred M'$,
  म्हणजे $M \ident M_1 \bredpar \dots \bredpar M_k \ident M'$.
  \ref{lam:cr:b:lem:par-red} नुसार $M \ident M_1 \bred
  \dots \bred M_k \ident M'$. $\bred$ संक्रमक
  असल्याने $M \bred M'$.
\end{proof}

\begin{thm}\label{lam:cr:b:thm:cr}
  $\bred$ ला चर्च–रॉसर गुणधर्म आहे.
\end{thm}

\begin{proof}
  \ref{lam:cr:dap:thm:str}, \ref{lam:cr:pb:thm:cr} आणि
  \ref{lam:cr:b:lem:str} वरून लगेच मिळते.
\end{proof}












\section{समांतर $\beta\eta$-न्यूनीकरण}\label{lam:cr:pbe:sec}

या विभागात $\beta\eta$-न्यूनीकरणाशी निगडित समांतर
$\beta\eta$-न्यूनीकरणाला चर्च–रॉसर गुणधर्म आहे हे सिद्ध करू.

\begin{defn}[समांतर $\beta\eta$-न्यूनीकरण, $\beredpar$] \label{lam:cr:pbe:defn:beredpar}
  पदांवरील \emph{समांतर $\beta\eta$-न्यूनीकरण}
  ($\beredpar$) पुढीलप्रमाणे विगमनाने परिभाषित केले जाते:
  \begin{enumerate}
    \item \label{lam:cr:pbe:defn:beredpar1} $x \beredpar x$.
    \item \label{lam:cr:pbe:defn:beredpar2} $N \beredpar N'$ असेल, तर $\lambd[x][N] \beredpar
      \lambd[x][N']$.
    \item \label{lam:cr:pbe:defn:beredpar3} $P \beredpar P'$ आणि $Q \beredpar Q'$ असतील, तर $PQ \beredpar
      P'Q'$.
    \item \label{lam:cr:pbe:defn:beredpar4} $N \beredpar N'$ आणि $Q \beredpar Q'$ असतील, तर
      $(\lambd[x][N])Q \beredpar \Subst{N'}{Q'}{x}$.
    \item \label{lam:cr:pbe:defn:beredpar5} $N \beredpar N'$ असेल, तर $\lambd[x][Nx]
      \beredpar N'$, मात्र $x \notin FV(N)$ असले पाहिजे.
  \end{enumerate}
\end{defn}

\begin{thm}\label{lam:cr:pbe:thm:refl}
  $M \beredpar M$.
\end{thm}

\begin{proof}
  सरावासाठी.
\end{proof}

\begin{prob}
  \ref{lam:cr:pbe:thm:refl} सिद्ध करा.
\end{prob}

\begin{defn}[$\beta\eta$-पूर्ण विकास]\label{lam:cr:pbe:defn:becd}
  $M$ चा \emph{$\beta\eta$-पूर्ण विकास}~$\becd{M}$
  पुढीलप्रमाणे परिभाषित केला जातो. पाचवी अट लागू
  होत असल्यास दुसऱ्या अटीऐवजी ती वापरायची आहे:
  \begin{align}
    \becd{x} &= x \label{lam:cr:pbe:defn:becd1} \\
    \becd{(\lambd[x][N])} &= \lambd[x][\becd{N}] \label{lam:cr:pbe:defn:becd2}\\
    \becd{(PQ)} &= \becd{P}\becd{Q} && \text{जर $P$ हे $\lambd$-अमूर्तीकरण नसेल तर}
    \label{lam:cr:pbe:defn:becd3} \\
    \becd{((\lambd[x][N])Q)} &= \Subst{\becd{N}}{\becd{Q}}{x}
                             \label{lam:cr:pbe:defn:becd4} \\
    \becd{(\lambd[x][Nx])} &= \becd{N} \label{lam:cr:pbe:defn:becd5} & \text{जर $x
                                                      \notin FV(N)$}
  \end{align}
\end{defn}

\begin{lem}\label{lam:cr:pbe:lem:comp}
  $M \beredpar M'$ आणि $R \beredpar R'$ असतील, तर $\Subst{M}{R}{y}
  \beredpar \Subst{M'}{R'}{y}$.
\end{lem}

\begin{proof}
  सर्व अमूर्तीकरणांची नावे आदेशनाच्या चरापासून आणि
  दोन्ही आदेशित पदांच्या मुक्त चरांपासून वेगळी निवडू.
  $M \beredpar M'$ च्या निष्पत्ती वर विगमन करू.

  पहिली चार प्रकरणे \ref{lam:cr:pb:lem:comp} मधील
  प्रकरणांसारखीच आहेत. अखेरचा नियम
  \ref{lam:cr:pbe:defn:beredpar5} असेल, तर काही $x$ आणि~$N'$
  साठी $M$ हे $\lambd[x][Nx]$, $M'$ हे $N'$
  आहे, $x \notin FV(N)$ आणि $N \beredpar N'$.
  आपल्याला $\Subst{(\lambd[x][Nx])}{R}{y} \beredpar
  \Subst{N'}{R'}{y}$, म्हणजे
  $\lambd[x][\Subst{N}{R}{y} x] \beredpar \Subst{N'}{R'}{y}$,
  हे सिद्ध करायचे आहे. ते \ref{lam:cr:pbe:defn:beredpar}\ref{lam:cr:pbe:defn:beredpar5}
  आणि विगमन गृहीतकावरून मिळते. बाह्य बद्ध नाव मूळ
  देहात मुक्त नाही आणि आदेशित पदातही मुक्त नाही, म्हणून
  आदेशनानंतरच्या देहात ते मुक्त होत नाही. त्यामुळे वापरलेल्या
  इटा नियमाची ताजेपणाची अट पूर्ण होते.
\end{proof}

\begin{lem}\label{lam:cr:pbe:lem:cont}
  $M \beredpar M'$ असेल, तर $M' \beredpar \becd{M}$.
\end{lem}

\begin{proof}

प्रथम समांतर न्यूनीकरणाने नवीन मुक्त चरे निर्माण होत नाहीत,
हे त्याच्या पाच नियमांवरील विगमनाने मिळते. बीटा प्रकरणात
आदेशनाचे मुक्त-चर नियम वापरतात; इटा प्रकरणात बद्ध चर
उरलेल्या पदात मुक्त नसतो. पूर्ण विकास हाही एक समांतर
न्यूनीकरण आहे, हे त्याच्या व्याख्येवरील विगमनाने मिळते.
म्हणून पूर्ण विकासानेही मुक्त चरे वाढत नाहीत. खाली आवश्यक
तेथे प्रतिनिधींची बद्ध नावे सर्व संबंधित सांत मुक्त-चर
संचांपासून वेगळी निवडू.
पदाच्या रचनेवर पुढील दोन दावे एकत्र सिद्ध करू:
\begin{enumerate}
\item मुख्य दावा: $M\beredpar M'$ असल्यास
$M'\beredpar\becd{M}$.
\item सहाय्यक दावा: $A=\lambd[x][B]$, $A\beredpar A'$
आणि $S\beredpar T$ असल्यास
$A'S\beredpar\Subst{\becd{B}}{T}{x}$.
\end{enumerate}
सहाय्यक दावा अमूर्तीकरणांपुरताच आहे. तो प्रथम पाहू.
$A\beredpar A'$ ची शेवटची पायरी अमूर्तीकरणाचा नियम
असेल, तर $A'=\lambd[x][B']$ आणि $B\beredpar B'$.
लहान पदासाठीच्या मुख्य विगमन दाव्याने
$B'\beredpar\becd{B}$ मिळते. याला $S\beredpar T$
जोडून बीटा समांतर नियमाने सहाय्यक निष्कर्ष मिळतो.
शेवटची पायरी इटा नियमाची असेल, तर $B=Hx$,
$x\notin\FV{H}$, $A'=H'$ आणि $H\beredpar H'$.
$H$ अमूर्तीकरण नसेल, तर लहान पदासाठीच्या मुख्य दाव्याने
$H'\beredpar\becd{H}$ मिळते. उपयोजनाच्या नियमाने
$H'S\beredpar\becd{H}T$; ही उजवी बाजू
$\Subst{\becd{(Hx)}}{T}{x}$ आहे.
$H=\lambd[y][L]$ असेल, तर लहान अमूर्तीकरणासाठीचा
सहाय्यक विगमन दावा
$H'S\beredpar\Subst{\becd{L}}{T}{y}$ देतो.
येथे $y$ हे $x$ पासून आणि $T$ च्या मुक्त चरांपासून
वेगळे निवडले आहे. पूर्ण विकासाच्या चौथ्या नियमाने
$\becd{(Hx)}=\Subst{\becd{L}}{x}{y}$; तसेच $x$
हे $L$ मध्ये किंवा त्याच्या पूर्ण विकासात मुक्त नाही.
\ref{lam:cr:pb:lem:comp} च्या सिद्धतेत दिलेल्या आदेशनसूत्राने
\[\Subst{\Subst{\becd{L}}{x}{y}}{T}{x}
=\Subst{\becd{L}}{T}{y}.
\]
म्हणून या प्रकरणातही सहाय्यक दावा सिद्ध होतो.

आता मुख्य दावा सिद्ध करू. चरासाठी तो परावर्तितेने मिळतो.
$M=PQ$ अमूर्तीकरण नसलेल्या फलनस्थानील पदासह असेल,
तर समांतर उपयोजनाचाच नियम लागू होतो. दोन्ही उपपदांना
मुख्य विगमन दावा लावून $P'Q'\beredpar\becd{P}\becd{Q}$
मिळते; उजवी बाजू $\becd{M}$ आहे.
$M=(\lambd[x][N])Q$ असेल आणि शेवटची पायरी बीटा
नियमाची असेल, तर $M'=\Subst{N'}{Q'}{x}$,
$N\beredpar N'$ आणि $Q\beredpar Q'$.
दोन्ही मुख्य विगमन दावे आणि \ref{lam:cr:pbe:lem:comp} यांनी
$M'\beredpar\Subst{\becd{N}}{\becd{Q}}{x}=\becd{M}$
मिळते. शेवटची पायरी उपयोजनाची असेल, तर
$M'=P'Q'$, $\lambd[x][N]\beredpar P'$ आणि
$Q\beredpar Q'$. मुख्य विगमन दाव्याने
$Q'\beredpar\becd{Q}$. मग लहान फलनस्थानील
अमूर्तीकरणासाठीचा सहाय्यक दावा तोच आवश्यक निष्कर्ष देतो;
यात $P'$ अजून अमूर्तीकरण असणे गृहीत धरलेले नाही.
शेवटी $M=\lambd[x][N]$ पाहू. शेवटची पायरी इटा
नियमाची असेल, तर $N=Hx$, $x\notin\FV{H}$ आणि
$M'=H'$ जेथे $H\beredpar H'$. मुख्य विगमन दाव्याने
$H'\beredpar\becd{H}=\becd{M}$.
शेवटची पायरी अमूर्तीकरणाच्या नियमाची असेल, तर
$M'=\lambd[x][N']$ जेथे $N\beredpar N'$.
$M$ इटा रेडेक्स नसेल, तर मुख्य विगमन दावा आणि
अमूर्तीकरणाचा नियम थेट निष्कर्ष देतात.
इटा रेडेक्स असेल, तर $N=Hx$ आणि $x\notin\FV{H}$.
$Hx\beredpar N'$ ची शेवटची पायरी उपयोजनाची असेल,
तर $N'=H'x$, $H\beredpar H'$; मुख्य विगमन दाव्याने
$H'\beredpar\becd{H}$ आणि इटा नियमाने
$\lambd[x][H'x]\beredpar\becd{H}=\becd{M}$.
ही पायरी बीटाची असेल, तर $H=\lambd[y][L]$,
$L\beredpar L'$ आणि $N'=\Subst{L'}{x}{y}$.
नवीन मुक्त चरे निर्माण होत नसल्याने $x\notin\FV{L'}$,
आणि $x\ne y$ अशी नावांची निवड आहे. म्हणून
$\lambd[x][\Subst{L'}{x}{y}]$ हे $\lambd[y][L']$
याच वर्गपदाचे दुसरे रूप आहे. अमूर्तीकरणाच्या नियमाने
$H\beredpar\lambd[y][L']$, आणि लहान पद $H$ साठीचा
मुख्य विगमन दावा
$\lambd[y][L']\beredpar\becd{H}=\becd{M}$ देतो.
ही सर्व शक्य प्रकरणे आहेत. प्रत्येक वापरलेल्या विगमन
दाव्याचे मूलपद सध्याच्या पदाचे योग्य उपपद आहे,
म्हणून एकत्र विगमन पूर्ण झाले.
\end{proof}

\begin{thm}\label{lam:cr:pbe:thm:cr}
  $\beredpar$ ला चर्च–रॉसर गुणधर्म आहे.
\end{thm}

\begin{proof}
  \ref{lam:cr:pbe:lem:cont} वरून लगेच मिळते.
\end{proof}












\section{$\beta\eta$-न्यूनीकरण}\label{lam:cr:be:sec}

$\beta\eta$-न्यूनीकरणाला ($\bered$) चर्च–रॉसर
गुणधर्म आहे.

\begin{lem}\label{lam:cr:be:lem:one-par}
  $M \beredone M'$ असेल, तर $M \beredpar M'$.
\end{lem}

\begin{proof}
  $M \beredone M'$ च्या निष्पत्ती वर विगमन करू.
  $M \eredone M'$ हे $\eta$-परिवर्तनाचे प्रकरण
  असेल (म्हणजे \ref{lam:syn:eta:defn:beredone}),
  तर समांतर इटा-नियमासाठी \ref{lam:cr:pbe:thm:refl}
  वापरू. इतर प्रकरणे \ref{lam:cr:b:lem:one-par}
  प्रमाणेच आहेत.
\end{proof}

\begin{lem}\label{lam:cr:be:lem:par-red}
  $M \beredpar M'$ असेल, तर $M \bered M'$.
\end{lem}

\begin{proof} $M \beredpar M'$ च्या निष्पत्ती वर विगमन करू.

  पहिली चार प्रकरणे समांतर बीटा-न्यूनीकरणाच्या
  संबंधित प्रकरणांप्रमाणे हाताळता येतात. अखेरचा नियम
  \ref{lam:cr:pbe:defn:beredpar5} असेल, तर काही $x$, $N$,
  $N'$ साठी $M$ हे $\lambd[x][Nx]$ आणि $M'$ हे $N'$
  आहे, जेथे $x \notin FV(N)$ आणि $N \beredpar N'$.
  म्हणून आधी $\lambd[x][Nx]$ चे $N$ मध्ये
  $\eta$-परिवर्तनाने न्यूनीकरण करू आणि मग विगमन
  गृहीतकानुसार $N \bered N'$ दाखवणाऱ्या $\beredone$
  पायऱ्या घेऊ.
\end{proof}

\begin{lem}\label{lam:cr:be:lem:str}
  $\beredpar$ समाविष्ट करणारा सर्वांत लहान संक्रमक
  संबंध $\bered$ आहे.
\end{lem}

\begin{proof}
  \ref{lam:cr:b:lem:str} प्रमाणे.
\end{proof}

\begin{thm}\label{lam:cr:be:thm:cr}
  $\bered$ ला चर्च–रॉसर गुणधर्म आहे.
\end{thm}

\begin{proof}
  \ref{lam:cr:dap:thm:str}, \ref{lam:cr:pbe:thm:cr} आणि
  \ref{lam:cr:be:lem:str} वरून मिळते.
\end{proof}


